We remark that since there is an extension operator from $H^k(\mathbb{R}_+^m)$ to $H^k(\mathbb{R}^m)$, we also have a trace operator which maps a function in $H^k(\mathbb{R}_+^m)$ to its boundary values in $H^{k-1/2}(\mathbb{R}^{m-1})$. By using a partition of unity argument, we can extend this result to domains with bounded boundary.

\begin{Theorem}
\label{thm:smooth-domain-trace-and-right-inverse}
Let $k$ be a positive integer. Assume that $\Omega$ is of class $C^k$ and $\partial\Omega$ is bounded. Then there is a bounded trace operator $\Gamma:H^k(\Omega)\to H^{k-1/2}(\partial\Omega)$. Moreover, $\Gamma$ has a bounded right inverse.
\end{Theorem}

If $l<k$, then the $l$th derivatives have traces in $H^{k-l-1/2}(\partial\Omega)$. It is customary to formulate trace theorems involving higher derivatives in terms of derivatives in the direction normal to $\partial\Omega$.

\begin{Theorem}
\label{thm:domain-higher-normal-trace}
Let $k,l$ be positive integers such that $k>l$. Let $\Omega$ be of class $C^k$ and let $\partial\Omega$ be bounded. Then there exists a continuous trace operator
\begin{equation}\label{eq:domain-higher-normal-trace-product-map}
T_l:H^k(\Omega)\to\prod_{j=0}^lH^{k-j-1/2}(\partial\Omega)
\end{equation}
with the property that
\begin{equation}\label{eq:domain-higher-normal-trace-smooth-values}
T\phi=\left(\phi,\dfrac{\partial\phi}{\partial n},\ldots,\dfrac{\partial^l\phi}{\partial n^l}\right)
\end{equation}
for every smooth $\phi$. The operator $T_l$ has a bounded right inverse.
\end{Theorem}

We can now characterize $H_0^k(\Omega)$ in terms of boundary conditions.

\begin{Theorem}
\label{thm:zero-boundary-sobolev-characterization}
Let $\Omega$ be of class $C^k$ and let $\partial\Omega$ be bounded. Then $H_0^k(\Omega)$ is the set of all those functions in $u\in H^k(\Omega)$ for which
\begin{equation}\label{eq:zero-boundary-sobolev-normal-trace-characterization}
u=\dfrac{\partial u}{\partial n}=\cdots=\dfrac{\partial^{k-1}u}{\partial n^{k-1}}=0
\end{equation}
on $\partial\Omega$ in the sense of trace.
\end{Theorem}

\begingroup
\postdisplaypenalty=10000
\renewcommand{\endProof}{\par\nobreak\hfill$\blacksquare$\par\vspace{0.1cm}}
\begin{Proof}
If $u\in\mathcal{D}(\Omega)$, it is clear that \eqref{eq:zero-boundary-sobolev-normal-trace-characterization} holds. By continuity, \eqref{eq:zero-boundary-sobolev-normal-trace-characterization} then holds for $u\in H_0^k(\Omega)$. We need to establish the converse. By using a partition of unity and local coordinate transformations, we are reduced to the case $\Omega=\mathbb{R}_+^m$. Let now $k=1$ and let $u\in H^1(\mathbb{R}_+^m)$ be such that $u(\bm{x}^{\prime},0)=0$ in the sense of trace. Let $Eu$ be the extension of $u$ by zero. To show that $Eu\in H^1(\mathbb{R}^m)$, it suffices to establish that $\partial(Eu)/\partial x_i=E(\partial u/\partial x_i)$. This is clear for $i<m$. For $i=m$, we have, for any $\phi\in\mathcal{D}(\mathbb{R}^m)$:
\begin{equation}\label{eq:zero-extension-normal-weak-derivative-calculation}
\begin{aligned}
\int_{\mathbb{R}^m}\dfrac{\partial\phi}{\partial x_m}(Eu)\,\mathrm{d}\bm{x}
&=\int_{\mathbb{R}_+^m}\dfrac{\partial\phi}{\partial x_m}u\,\mathrm{d}\bm{x}\\
&=\int_{\mathbb{R}^{m-1}}\phi(\bm{x}^{\prime},0)u(\bm{x}^{\prime},0)\,\mathrm{d}\bm{x}^{\prime}-\int_{\mathbb{R}_+^m}\phi\dfrac{\partial u}{\partial x_m}\,\mathrm{d}\bm{x}\\
&=-\int_{\mathbb{R}^m}\phi E\dfrac{\partial u}{\partial x_m}\,\mathrm{d}\bm{x}.
\end{aligned}
\end{equation}
An analogous argument applies to higher derivatives. Once we know that $Eu\in H^k(\mathbb{R}^m)$, the rest follows by considering the sequence $u_n=Eu(\bm{x}-\dfrac{1}{n}\bm{e}_m)$. Since the support of $u_n$ is bounded away from $\partial\mathbb{R}_+^m$, it is easy to approximate $u_n$ by test functions.
\end{Proof}
\endgroup
